% ============================================================
% Chapter 4: Proposed Method: Specification and Design
%   Section: The Fact-Checking Pipeline
%     Subsection: Admission Filter
% Included from chapters/04-proposed-method/06-fact-checking-pipeline/fact-checking-pipeline.tex
% Sources: docs/arquitectura-tecnica.md 2.4.1;
% backend/src/services/admission/README.md.
% ============================================================

\subsection{Admission Filter}
\label{subsec:admission}

Before anything is spent on verification, the article must pass three
independent checks. They judge the \emph{article} --- its topics, its
tone, its neighbours in the vector store --- and none needs evidence, a
search or a language model, which is why they come first and are the
cheapest stage.

\begin{itemize}
    \item \textbf{Topic relevance}: at least one topic's confidence
    must reach a minimum (0.35 by default). No topic has to dominate,
    since an article can legitimately touch several.
    \item \textbf{Positive impact}: a weighted average of
    constructiveness, inspirational value, hopefulness, positive
    sentiment, objectivity and social impact, in $[0,1]$, must reach a
    minimum (0.1 by default). Hard floors on objectivity and negative
    sentiment, which would reject an article whatever its average, exist
    behind a single switch and are off by default: the floors on
    constructiveness and inspiration rejected too many real articles,
    and now only lower the score and are recorded as reasons.
    \item \textbf{Not a duplicate}: the article's embedding is compared
    with its nearest stored neighbour. Below a relatedness threshold
    (0.80) the two are unrelated; between it and a duplicate threshold
    (0.90) they are \emph{related}, a useful signal that does not
    disqualify; above it the article is a duplicate. A stored article
    with the same URL never counts: it is the same article from an
    earlier run, and a re-analysis must not reject itself.
\end{itemize}

All three always run, even after one fails, so a rejection names every
reason. A rejected article is stored with its reasons and is never
searched or sent to a model. Every threshold can be overridden for one
run (Section~\ref{sec:dev-storage}).
